% ============================================================
%  Κεφάλαιο 2 — Βιβλιογραφική Ανασκόπηση
% ============================================================
\chapter{Βιβλιογραφική Ανασκόπηση}
\label{ch:literature}

% ============================================================
\section{Πρόβλεψη Τιμών Ηλεκτρικής Ενέργειας}
\label{sec:epf_review}
% ============================================================

\subsection{Εξέλιξη πεδίου και βασικές ανασκοπήσεις}

Η βραχυπρόθεσμη πρόβλεψη τιμών ηλεκτρικής ενέργειας
(\textit{Electricity Price Forecasting}, EPF) είναι ένα από τα πιο απαιτητικά προβλήματα χρονοσειρών~\cite{weron2014electricity}.
Ο Weron~\cite{weron2014electricity} παρουσιάζει εκτεταμένη ανασκόπηση,
τονίζοντας ότι οι αγορές ηλεκτρισμού παρουσιάζουν δύσκολη
δυναμική: απότομες αιχμές (\textit{spikes}), ισχυρή εποχικότητα
(ωριαία, ημερήσια, εβδομαδιαία) και μη σταθερή διακύμανση.
Οργανώνει τις μεθόδους σε πέντε κατηγορίες: multi-agent,
fundamental, reduced-form, statistical και computational intelligence,
τεκμηριώνοντας ότι οι στατιστικές και υπολογιστικές προσεγγίσεις
κυριαρχούν στη βραχυπρόθεσμη πρόβλεψη.

Οι Lago et al.~\cite{lago2021} επικεντρώνονται σε αλγόριθμους
τελευταίας γενιάς και προτείνουν ένα ανοιχτό σημείο αναφοράς (benchmark) για
αντικειμενική σύγκριση μεθόδων στην αγορά Nord Pool.
Ο Kanousis~\cite{kanousis2025} εφαρμόζει ένα σύνολο μεθόδων ML
στην ελληνική αγορά επόμενης ημέρας (\textit{Day-Ahead Market}),
επιτυγχάνοντας σημαντική βελτίωση έναντι στατιστικών μοντέλων αναφοράς.
Ο Κοντράρος~\cite{kontraros2023} εξετάζει τεχνικές ensemble
για ελληνικά δεδομένα, αναδεικνύοντας τη σημασία της επιλογής
χαρακτηριστικών στα αποτελέσματα πρόβλεψης.

\subsection{Κατηγορίες Μεθόδων Πρόβλεψης}

Τα μοντέλα EPF διακρίνονται σε τέσσερις ευρείες κατηγορίες.
Τα \emph{στατιστικά μοντέλα} (ARIMA, GARCH, VAR) κυριαρχούσαν
στη βιβλιογραφία έως τα μέσα του 2010.
Τα \emph{μοντέλα μηχανικής μάθησης} (Gradient Boosting Machines, SVR, MLP)
συχνά υπερτερούν των στατιστικών μοντέλων σε σύγχρονες αγορές,
κατά κύριο λόγο λόγω της ικανότητάς τους να ενσωματώνουν
πολλαπλές εξωγενείς μεταβλητές χωρίς ισχυρές παραδοχές κατανομής.
Τα \emph{μοντέλα βαθιάς μάθησης} (LSTM, Transformer)
επιτυγχάνουν ισχυρά αποτελέσματα σε παγκόσμιες αγορές υψηλής
ρευστότητας~\cite{lago2018,marcjasz2023beating,lim2021temporal}. Ειδικότερα, το
\textit{Temporal Fusion Transformer}~\cite{lim2021temporal}
συνδυάζει μηχανισμούς attention (προσοχής, που σταθμίζουν δυναμικά τη βαρύτητα διαφορετικών χρονικών βημάτων) με temporal self-attention, επιτρέποντας
ερμηνεύσιμη μοντελοποίηση πολλαπλών χρονικών οριζόντων.
Ωστόσο, οι Zeng et al.~\cite{zeng2023transformers} υποστηρίζουν ότι απλά
γραμμικά μοντέλα συχνά ξεπερνούν τους Transformers σε χρονοσειρές
περιορισμένου μήκους, όπως η ωριαία EPF ($T \approx 80.000$ δείγματα).
Για την ελληνική αγορά ειδικότερα, οι μέθοδοι ML (Gradient Boosting)
παραμένουν ανταγωνιστικές λόγω περιορισμένου μεγέθους δεδομένων
και υψηλής μεταβλητότητας που δεν ευνοεί την πλήρη αξιοποίηση deep learning.
Τέλος, τα \emph{υβριδικά/ensemble μοντέλα} συνδυάζουν προβλέψεις
πολλαπλών αλγόριθμων με τεχνικές σταθμισμένης συγχώνευσης,
συχνά επιτυγχάνοντας καλύτερη γενίκευση από κάθε μεμονωμένο μοντέλο~\cite{dietterich2000,nowotarski2014}.

\subsection{Ο ρόλος εποχικής διόρθωσης}

Η εποχική αποσύνθεση πριν από την εφαρμογή ML έχει αναδειχθεί ως χρήσιμη
προεπεξεργασία για την EPF.
Οι Uniejewski et al.~\cite{uniejewski2019} δείχνουν ότι η διόρθωση
βάσει της μακροπρόθεσμης (ετήσιας) εποχικής συνιστώσας (LTSC)
μειώνει σημαντικά το MAE σε ευρωπαϊκές αγορές. Οι Marcjasz et
al.~\cite{marcjasz2019eneco} επεκτείνουν αυτή την ιδέα σε
τεχνικές ειδικές ανά ώρα, τεκμηριώνοντας σημαντικές βελτιώσεις
σε δείκτες ακρίβειας. Στην παρούσα εργασία, οι εποχικές
ιδιότητες ενσωματώνονται απευθείας ως χαρακτηριστικά υστέρησης και κυκλικές κωδικοποιήσεις χρόνου (lag features / cyclic encodings) αντί ρητής αποσύνθεσης, ακολουθώντας
τη βέλτιστη πρακτική που περιγράφεται στον Lago et al.~\cite{lago2021}.


% ============================================================
\section{Πρόβλεψη Ηλεκτρικού Φορτίου}
\label{sec:elf_review}
% ============================================================

\subsection{Ιστορική εξέλιξη}

Η πρόβλεψη ηλεκτρικού φορτίου (\textit{Electricity Load
Forecasting}, ELF) έχει μακρύτερη ιστορία από την EPF.
Οι Gross και Galiana~\cite{gross1987} τεκμηριώνουν τις
πρώτες συστηματικές προσεγγίσεις ήδη από τη δεκαετία
του 1980, βασισμένες κυρίως σε παλινδρόμηση και εκθετική
εξομάλυνση. Ο Μαγκλαράς~\cite{magklaras2018} εξετάζει
τις σύγχρονες τεχνικές ML για το ελληνικό δίκτυο,
τονίζοντας τη σημασία των μετεωρολογικών χαρακτηριστικών.

\subsection{Κατηγοριοποίηση κατά Ορίζοντα}

Η βιβλιογραφία διακρίνει τέσσερις κατηγορίες ανάλογα
με τον ορίζοντα πρόβλεψης:

\begin{itemize}
  \item \textbf{VSTLF} (\textit{Very Short-Term Load Forecasting}):
        μερικά λεπτά έως μία ώρα, κρίσιμο για τον έλεγχο
        ισορροπίας σε πραγματικό χρόνο.
  \item \textbf{STLF} (\textit{Short-Term Load Forecasting}):
        μία ώρα έως δύο εβδομάδες, ο ορίζοντας της παρούσας εργασίας ($H=24$h).
  \item \textbf{MTLF} (\textit{Medium-Term Load Forecasting}):
        δύο εβδομάδες έως τρία χρόνια, σχεδιασμός συντήρησης δικτύου.
  \item \textbf{LTLF} (\textit{Long-Term Load Forecasting}):
        πάνω από τρία χρόνια, επενδυτικές αποφάσεις υποδομής.
\end{itemize}

\subsection{Στατιστικά Χαρακτηριστικά}

Το ηλεκτρικό φορτίο παρουσιάζει τρία διακριτά
στατιστικά χαρακτηριστικά. Πρώτον, ισχυρή πολλαπλή
\emph{εποχικότητα}: ημερήσιο μοτίβο (υψηλά φορτία εργάσιμες
ώρες), εβδομαδιαίο μοτίβο (χαμηλά Σαββατοκύριακα) και ετήσιο
μοτίβο (αιχμές χειμώνα/καλοκαιριού λόγω θέρμανσης/ψύξης).
Δεύτερον, ισχυρή \emph{αυτοσυσχέτιση} στις υστερήσεις
$t-1$ (1 ώρα, η προηγούμενη ώρα), $t-24$ (η ίδια ώρα χθες) και
$t-168$ (η ίδια ώρα την ίδια εργάσιμη/μη-εργάσιμη μέρα
την προηγούμενη εβδομάδα). Τρίτον, \emph{μη γραμμικότητα}
λόγω ακραίων κλιματικών συνθηκών (θερμοκρασία, υγρασία),
ειδικών ημερών (εορτές, απεργίες) και απότομων αλλαγών
συμπεριφοράς του καταναλωτή.

Αυτά τα χαρακτηριστικά καθιστούν τα μοντέλα ML
κατάλληλα, καθώς μπορούν να συλλάβουν
μη γραμμικές αλληλεπιδράσεις μεταξύ πολλαπλών
χαρακτηριστικών χωρίς ρητή μοντελοποίηση.


% ============================================================
\section{Μοντέλα Αναφοράς (Naive Baselines)}
\label{sec:baselines_review}
% ============================================================

Πριν από κάθε αλγόριθμο μηχανικής μάθησης, η βιβλιογραφία πρόβλεψης χρονοσειρών ορίζει
ένα σύνολο αφελών μοντέλων αναφοράς (naive baselines): μεθόδους εμμονής (persistence)
χωρίς εκπαίδευση, που θέτουν το ελάχιστο όριο επίδοσης το οποίο οφείλει να υπερβεί κάθε
μη τετριμμένο μοντέλο~\cite{lago2021}. Στην εργασία χρησιμοποιούνται τέσσερα: το
\textbf{Naive-1} (τιμή προηγούμενης ώρας), το \textbf{Naive-24} (ίδια ώρα προηγούμενης
ημέρας), το \textbf{Naive-168} (ίδια ώρα προηγούμενης εβδομάδας) και το \textbf{Seasonal
Profile} (μέσος όρος ανά ώρα και ημέρα εβδομάδας του συνόλου εκπαίδευσης). Τα τρία πρώτα
είναι η άμεση λειτουργική έκφραση της αυτοσυσχέτισης που περιγράφηκε παραπάνω (υστερήσεις
$t-1$, $t-24$, $t-168$). Η επιλογή του κατάλληλου baseline εξαρτάται από το σύνολο
πληροφορίας κάθε στρατηγικής: το Naive-1 για τις μονοβηματικές προσεγγίσεις, το Naive-24
για τις στρατηγικές μίας αφετηρίας (MIMO/Direct), ζήτημα στο οποίο επανερχόμαστε στην
Ενότητα~\ref{sec:multistep}. Οι τυπικοί ορισμοί δίνονται στο
Κεφάλαιο~\ref{ch:methods}.


% ============================================================
\section{Αλγόριθμοι Μηχανικής Μάθησης}
\label{sec:ml_algorithms}
% ============================================================

\subsection{Μέθοδοι Gradient Boosting}
\label{subsec:gbt}

\paragraph{Random Forest.}
Ο αλγόριθμος Random Forest~\cite{breiman2001rf} εκπαιδεύει
$B$ ανεξάρτητα δέντρα απόφασης, το καθένα σε bootstrap
δείγμα του συνόλου εκπαίδευσης, και συνδυάζει τις
προβλέψεις τους με αριθμητικό μέσο:
%
\begin{equation}
  \hat{y} = \frac{1}{B}\sum_{b=1}^{B} T_b(x)
  \label{eq:rf}
\end{equation}
%
Η διάσπαση κάθε κόμβου επιλέγεται ώστε να μεγιστοποιεί τη μείωση της
διακύμανσης (\textit{variance reduction}):
%
\begin{equation}
  \mathrm{IV} = \frac{1}{N}\sum_{i \in I_L}\left(y_i-\bar{y}_L\right)^2
              + \frac{1}{N}\sum_{i \in I_R}\left(y_i-\bar{y}_R\right)^2
  \label{eq:rf_var}
\end{equation}
%
όπου $I_L, I_R$ τα δείγματα αριστερού/δεξιού κόμβου και $\bar{y}_L, \bar{y}_R$
οι αντίστοιχοι μέσοι όροι.

Ως μέθοδος \emph{ομαδικής μάθησης} (\textit{ensemble learning}), το Random Forest ακολουθεί την τεχνική \emph{bagging} (Bootstrap Aggregating): αντί να εμπιστεύεται ένα μόνο δέντρο, συνδυάζει πολλά ανεξάρτητα και λειτουργεί σαν επιτροπή εκτιμητών, όπου ο μέσος όρος πολλών κρίσεων είναι πιο σταθερός και αξιόπιστος από μία μεμονωμένη.

Η τυχαία επιλογή υποσυνόλου χαρακτηριστικών σε κάθε
κόμβο μειώνει τη διακύμανση χωρίς αντίστοιχη αύξηση
του σφάλματος μεροληψίας. Η μέθοδος είναι ιδιαίτερα ανθεκτική σε ακραίες τιμές (outliers) και σε υπερπροσαρμογή (\textit{overfitting})~\cite{breiman2001rf}.

\begin{figure}[htbp]
  \centering
  \includegraphics[width=0.95\textwidth]{fig_random_forest}
  \caption{Σχηματική αναπαράσταση του αλγόριθμου Random Forest:
    $B$ δέντρα εκπαιδεύονται σε bootstrap δείγματα και η τελική
    πρόβλεψη προκύπτει ως ο μέσος όρος τους.}
  \label{fig:rf}
\end{figure}

\paragraph{XGBoost.}
Ο αλγόριθμος XGBoost (\textit{Extreme Gradient Boosting})~\cite{chen2016xgboost}
κατασκευάζει δέντρα σειριακά, με κάθε νέο δέντρο να
μαθαίνει τα κατάλοιπα του προηγούμενου. Η αντικειμενική
συνάρτηση εκπαίδευσης περιλαμβάνει ρητό όρο κανονικοποίησης:
%
\begin{equation}
  \mathcal{L} = \sum_{i} \ell\!\left(y_i,\,\hat{y}_i\right)
    + \sum_{k}\Omega(f_k)
  \label{eq:xgb}
\end{equation}
%
όπου $\ell$ είναι η συνάρτηση απώλειας (π.χ.\ τετραγωνικό σφάλμα)
και $\Omega(f_k) = \gamma T_k + \frac{1}{2}\lambda\|w_k\|^2$
ελέγχει την πολυπλοκότητα κάθε δέντρου $f_k$ μέσω
του αριθμού φύλλων $T_k$ και των βαρών $w_k$.
Η πρόβλεψη ενημερώνεται προσθετικά,
$\hat{y}_i^{(t)} = \hat{y}_i^{(t-1)} + f_t(x_i)$, και η αντικειμενική συνάρτηση
στο βήμα $t$ προσεγγίζεται με ανάπτυγμα Taylor δεύτερης τάξης:
%
\begin{equation}
  \mathcal{L}^{(t)} \approx \sum_{i=1}^{N}\left[g_i f_t(x_i)
     + \tfrac{1}{2}h_i f_t^2(x_i)\right]
     + \gamma T + \tfrac{1}{2}\lambda\sum_{j=1}^{T} w_j^2
  \label{eq:xgb_taylor}
\end{equation}
%
όπου $g_i$ και $h_i$ είναι η πρώτη και η δεύτερη παράγωγος της απώλειας ως προς
την πρόβλεψη, $T$ ο αριθμός φύλλων και $w_j$ τα βάρη τους.
Ο αλγόριθμος χρησιμοποιεί \emph{level-wise} ανάπτυξη
δέντρου, δηλαδή διαχωρίζει όλους τους κόμβους του
ίδιου επιπέδου πριν προχωρήσει στο επόμενο.

\paragraph{LightGBM.}
Ο αλγόριθμος LightGBM~\cite{ke2017lgbm} βελτιώνει
τον XGBoost με δύο καινοτομίες. Πρώτον, η \emph{leaf-wise}
ανάπτυξη δέντρου επιλέγει πάντα το φύλλο με τη μεγαλύτερη
μείωση απώλειας, ανεξαρτήτως επιπέδου, επιτρέποντας
ασύμμετρα δέντρα που συχνά καλύπτουν καλύτερα τις μη-γραμμικές
σχέσεις. Δεύτερον, η τεχνική \textit{Gradient-based One-Side
Sampling} (GOSS) εστιάζει την εκπαίδευση στα δείγματα
με μεγαλύτερα gradients, μειώνοντας τον χρόνο υπολογισμού
χωρίς ουσιαστική απώλεια ακρίβειας.
Το κέρδος ενός υποψήφιου σημείου διαχωρισμού υπολογίζεται από τα αθροίσματα των
gradients ($G$) και hessians ($H$) στον αριστερό και δεξιό κόμβο:
%
\begin{equation}
  \mathrm{Gain} = \frac{1}{2}\left[
    \frac{G_L^2}{H_L+\lambda} + \frac{G_R^2}{H_R+\lambda}
    - \frac{(G_L+G_R)^2}{H_L+H_R+\lambda}\right]
  \label{eq:lgbm_gain}
\end{equation}
%
όπου $\lambda$ είναι ο συντελεστής κανονικοποίησης $L_2$. Για απώλεια L2
(\texttt{objective="regression"}), ισχύει $g_i=\hat{y}_i-y_i$ και $h_i=1$.


\begin{figure}[htbp]
  \centering
  \includegraphics[width=\textwidth]{fig_leafwise_levelwise}
  \caption{Σύγκριση \textit{level-wise} (XGBoost) και \textit{leaf-wise}
    (LightGBM) στρατηγικής ανάπτυξης δέντρου.}
  \label{fig:leafwise_levelwise}
\end{figure}

\subsection{Support Vector Regression}
\label{subsec:svr}

Η \textit{Support Vector Regression} (SVR)~\cite{smola2004}
επιδιώκει να βρει μια συνάρτηση $f(x)$ που
αποκλίνει από τις πραγματικές τιμές $y_i$ κατά
το πολύ $\varepsilon$, ελαχιστοποιώντας παράλληλα
την πολυπλοκότητα του μοντέλου:
%
\begin{equation}
  \min_{w,b,\xi,\xi^*}
    \;\frac{1}{2}\|w\|^2
    + C\sum_{i=1}^{n}\left(\xi_i + \xi_i^*\right)
  \label{eq:svr}
\end{equation}
%
υπό τους περιορισμούς $y_i - \langle w,x_i\rangle - b \leq \varepsilon + \xi_i$
και $\langle w,x_i\rangle + b - y_i \leq \varepsilon + \xi_i^*$,
με $\xi_i,\xi_i^* \geq 0$.
Η παράμετρος $C$ ελέγχει την ισορροπία μεταξύ
απλότητας μοντέλου και ανοχής σε σφάλμα,
ενώ η $\varepsilon$ ορίζει τον «σωλήνα ανοχής»
εντός του οποίου τα σφάλματα αγνοούνται.
Η ζώνη αυτή ορίζεται μέσω της $\varepsilon$-insensitive συνάρτησης απώλειας,
$L_\varepsilon(y_i,\hat{y}_i) = \max\!\left(0,\,|y_i-\hat{y}_i|-\varepsilon\right)$.

Η χρήση πυρήνα RBF δίνει τη δυνατότητα στο μοντέλο να
συλλάβει μη-γραμμικές σχέσεις μέσω της απεικόνισης
σε υψηλοδιάστατο χώρο χαρακτηριστικών.
%
\begin{equation}
  K(x_i,x_j) = \exp\!\left(-\gamma\,\|x_i-x_j\|^2\right)
  \label{eq:svr_rbf}
\end{equation}


\begin{figure}[htbp]
  \centering
  \includegraphics[width=0.92\textwidth]{fig_svr_etube}
  \caption{Ο «σωλήνας ανοχής» ($\varepsilon$-tube) της SVR.
    Τα δείγματα εντός του σωλήνα δεν συνεισφέρουν στην
    αντικειμενική συνάρτηση, μόνο τα δείγματα εκτός
    (μεταβλητές χαλαρότητας $\xi_i, \xi_i^*$) τιμωρούνται.}
  \label{fig:svr}
\end{figure}

\subsection{Multi-Layer Perceptron (MLP)}
\label{subsec:mlp}

Το \textit{Multi-Layer Perceptron} (MLP)~\cite{goodfellow2016deep}
υπολογίζει τη διέξοδό του μέσω διαδοχικών μετασχηματισμών
γραμμικής προβολής και μη-γραμμικής ενεργοποίησης:
%
\begin{equation}
  h^{(l)} = \sigma\!\left(
    W^{(l)}h^{(l-1)} + b^{(l)}
  \right)
  \label{eq:mlp}
\end{equation}
%
όπου $W^{(l)}$ και $b^{(l)}$ είναι
τα βάρη και τα όρια (\textit{biases}) του στρώματος $l$,
και $\sigma$ είναι η συνάρτηση ενεργοποίησης (ReLU
στη συγκεκριμένη υλοποίηση). Η ReLU ορίζεται ως
$\sigma(x) = \max(0, x)$, αποτρέποντας το πρόβλημα της εξαφάνισης κλίσης
(\textit{vanishing gradient}). Η εκπαίδευση γίνεται με
τον αλγόριθμο Adam, ο οποίος συνδυάζει προσαρμοστικό
βήμα εκμάθησης με ορμή δεύτερης τάξης. Κατά την εκπαίδευση χρησιμοποιούνται επιπλέον δύο τεχνικές κανονικοποίησης: το \emph{dropout}, που απενεργοποιεί τυχαία ένα ποσοστό νευρώνων σε κάθε επανάληψη ώστε το δίκτυο να μη στηρίζεται υπερβολικά σε μεμονωμένους νευρώνες (περιορισμός υπερπροσαρμογής), και η \emph{batch normalization}, που κανονικοποιεί τις ενεργοποιήσεις κάθε mini-batch, σταθεροποιώντας και επιταχύνοντας την εκπαίδευση.

\begin{figure}[htbp]
  \centering
  \includegraphics[width=0.92\textwidth]{fig_mlp_architecture}
  \caption{Αρχιτεκτονική MLP με κρυφά στρώματα.
    Κάθε κόμβος στρώματος $l$ δέχεται τις εξόδους
    του στρώματος $l-1$ και εφαρμόζει τον μετασχηματισμό
    της εξίσωσης~\ref{eq:mlp}.}
  \label{fig:mlp}
\end{figure}

Τα \emph{LSTM} (\textit{Long Short-Term Memory})~\cite{hochreiter1997lstm} είναι επανατροφοδοτούμενη επέκταση των MLP με ενσωματωμένους
μηχανισμούς «πύλης» για τη διαχείριση μακροπρόθεσμης εξάρτησης.
Στην παρούσα εργασία, τα LSTM δεν χρησιμοποιήθηκαν στο τελικό
σύνολο μοντέλων, καθώς τα δεδομένα ωριαίας ανάλυσης
και η ρητή συμπερίληψη lag χαρακτηριστικών ήδη
παρέχουν τη χρονική πληροφορία που αποσκοπεί
να εκμεταλλευτεί το LSTM.
Αντίστοιχα, οι αρχιτεκτονικές Transformer~\cite{vaswani2017attention}
έχουν επιδείξει ισχυρές επιδόσεις σε προβλήματα ακολουθιών, αλλά απαιτούν
μεγαλύτερα σύνολα δεδομένων από αυτά που διαθέτει η ελληνική αγορά κατά την
περίοδο μελέτης. Επιπλέον, πρόσφατη έρευνα~\cite{zeng2023transformers}
αμφισβητεί την υπεροχή τους έναντι απλούστερων μοντέλων σε χρονοσειρές
περιορισμένου μήκους.


% ============================================================
\section{Στρατηγικές Πολυβηματικής Πρόβλεψης}
\label{sec:multistep}
% ============================================================

Η πρόβλεψη σε ορίζοντα $H > 1$ βημάτων υλοποιείται
με τέσσερις στρατηγικές~\cite{taieb2012review,bontempi2012strategies}.

\paragraph{Teacher-Forcing (TF).}
Στη στρατηγική Teacher-Forcing (TF), το μοντέλο τροφοδοτείται με τις \emph{πραγματικές} υστερήσεις (lags) σε κάθε βήμα πρόβλεψης. Έχει νόημα μόνο σε συνθήκες αξιολόγησης ή σε σενάρια όπου οι
πραγματικές τιμές είναι διαθέσιμες (π.χ.\ ενδοημερήσια αγορά, online
ενημέρωση μοντέλου). Παρέχει το άνω φράγμα επίδοσης (oracle).

\paragraph{Recursive (Rec).}
Η αναδρομική στρατηγική εκπαιδεύει ένα μοντέλο
μονού βήματος και χρησιμοποιεί επαναληπτικά
τις \emph{προβλέψεις} ως εισόδους για τα επόμενα βήματα:
%
\begin{equation}
  \hat{y}_{t+h} = f\!\left(\hat{y}_{t+h-1},\,\ldots,\,\hat{y}_{t+1},\,y_t,\,y_{t-1},\,\ldots\right)
  \label{eq:recursive}
\end{equation}
%
Η μέθοδος είναι απλή στην υλοποίηση αλλά
υπόκειται σε συσσώρευση σφάλματος
(\textit{error accumulation}) για μεγάλους ορίζοντες.

\paragraph{MIMO.}
Η στρατηγική \textit{Multiple-Input Multiple-Output} (MIMO)
εκπαιδεύει ένα μοντέλο που παράγει ολόκληρο
το διάνυσμα πρόβλεψης ταυτόχρονα:
%
\begin{equation}
  \left(\hat{y}_{t+1},\,\ldots,\,\hat{y}_{t+H}\right)
    = f\!\left(x_t\right)
  \label{eq:mimo}
\end{equation}
%
Δεν υπάρχει συσσώρευση σφάλματος, αλλά η κοινή
εκπαίδευση για όλα τα βήματα μπορεί να οδηγήσει
σε μειωμένη ακρίβεια για τα πρώτα βήματα
του ορίζοντα σε σχέση με μεθόδους εξειδικευμένες ανά βήμα.

\paragraph{Direct.}
Η άμεση στρατηγική εκπαιδεύει $H$ ανεξάρτητα μοντέλα,
ένα για κάθε βήμα $h$ του ορίζοντα:
%
\begin{equation}
  \hat{y}_{t+h} = f_h\!\left(x_t\right),
    \quad h = 1,\,\ldots,\,H
  \label{eq:direct}
\end{equation}
%
Κάθε μοντέλο $f_h$ βελτιστοποιείται αποκλειστικά
για το συγκεκριμένο βήμα πρόβλεψης, εξαλείφοντας
τη συσσώρευση σφάλματος. Το κόστος είναι
$H$-πλάσιος χρόνος εκπαίδευσης.

\begin{figure}[htbp]
  \centering
  \includegraphics[width=\textwidth]{fig_multistep_strategies}
  \caption{Στρατηγικές πολυβηματικής πρόβλεψης. \textbf{Αριστερά:} η γενική
    ταξινόμηση της βιβλιογραφίας~\cite{taieb2012review}, με τις βασικές
    στρατηγικές (Recursive, Direct, MIMO) και τις υβριδικές παραλλαγές
    (DirRec, DIRMO). \textbf{Δεξιά:} οι τέσσερις στρατηγικές που
    αξιολογούνται στην παρούσα εργασία: \textbf{Teacher-Forcing (TF)},
    \textbf{Recursive (Rec)}, \textbf{Direct} και \textbf{MIMO}.}
  \label{fig:strategies}
\end{figure}

Στην παρούσα εργασία αξιολογούνται τέσσερις στρατηγικές
για τις ελληνικές αγορές τιμής και φορτίου, με ορίζοντα
$H = 24$ ώρες (ημερήσιος κύκλος DAM). Η στρατηγική TF
δεν εμφανίζεται ρητά στο σχήμα της βιβλιογραφίας γιατί
θεωρείται ειδική περίπτωση της Recursive που χρησιμοποιεί
\textit{πραγματικές} υστερήσεις αντί προβλέψεων κατά
την αξιολόγηση.

\paragraph{Δικαιοσύνη σύγκρισης και το «κόστος του ορίζοντα».}
Οι τέσσερις στρατηγικές δεν διαθέτουν το ίδιο σύνολο πληροφορίας τη στιγμή της
πρόβλεψης, γεγονός που καθιστά παραπλανητική την άμεση αντιπαραβολή των απόλυτων
σφαλμάτων τους. Το Teacher-Forcing «βλέπει» την πραγματική τιμή της αμέσως
προηγούμενης ώρας σε κάθε βήμα, ενώ τα MIMO και Direct προβλέπουν ολόκληρο το
24ωρο από μία και μόνη αφετηρία (το τέλος της προηγούμενης ημέρας), χωρίς πρόσβαση
σε ενδοημερήσιες πραγματικές τιμές. Αυτή η ασυμμετρία (το «κόστος του ορίζοντα») σημαίνει ότι ένα σφάλμα MIMO δεν είναι άμεσα συγκρίσιμο με ένα σφάλμα TF. Η
\emph{δίκαιη} αναφορά για τις στρατηγικές μίας αφετηρίας (MIMO/Direct) είναι το
\textbf{Naive-24}, η τιμή της ίδιας ώρας την προηγούμενη ημέρα, που μοιράζεται
ακριβώς το ίδιο σύνολο πληροφορίας. Ένα μοντέλο MIMO αποδεικνύει προστιθέμενη αξία
όταν ξεπερνά το Naive-24, και όχι όταν συγκρίνεται με το oracle TF. Η παρούσα
εργασία υιοθετεί αυτή την αρχή κατά την ερμηνεία των αποτελεσμάτων
(Κεφάλαιο~\ref{ch:experiments}).


% ============================================================
\section{Scheduled Sampling}
\label{sec:scheduled_sampling}
% ============================================================

Η τεχνική Scheduled Sampling (SS) προτάθηκε αρχικά
για αναδρομικά νευρωνικά δίκτυα~\cite{bengio2015scheduled}
ως ενδιάμεση λύση μεταξύ Teacher-Forcing (εκπαίδευση) και Recursive (συμπέρασμα). Το πρόβλημα που αντιμετωπίζει είναι η \emph{έκθεση σε μεροληψία} (\textit{exposure bias}): κατά την εκπαίδευση το μοντέλο βλέπει πάντα τις πραγματικές προηγούμενες τιμές, ενώ κατά την εκτέλεση τροφοδοτείται με τις δικές του προβλέψεις. Έτσι τα πρώιμα σφάλματα συσσωρεύονται και η ακολουθία αποκλίνει σταδιακά~\cite{bengio2015scheduled}. Ως στρατηγική \emph{σταδιακής μάθησης} (\textit{curriculum learning}), η μέθοδος μεταβαίνει ομαλά από την πλήρως καθοδηγούμενη εκπαίδευση με πραγματικές τιμές προς ένα σχήμα που στηρίζεται κυρίως στις παραγόμενες τιμές του ίδιου του μοντέλου. Κατά τη διάρκεια εκπαίδευσης,
κάθε εισερχόμενη υστέρηση επιλέγεται τυχαία:
%
\begin{equation}
  x_t = \begin{cases}
    y_{t-1}       & \text{με πιθανότητα } \varepsilon_t \\
    \hat{y}_{t-1} & \text{με πιθανότητα } 1 - \varepsilon_t
  \end{cases}
  \label{eq:ss}
\end{equation}
%
Η παράμετρος $\varepsilon_t$ μειώνεται μονότονα
κατά την εκπαίδευση (\textit{annealing}): ξεκινά
από $\varepsilon_0 = 1$ (καθαρό Teacher-Forcing)
και φτάνει στο $\varepsilon_T \approx 0$
(πλήρης αναδρομική πρόβλεψη). Στην αρχική διατύπωση, το πρόγραμμα μείωσης του $\varepsilon$ μπορεί να είναι γραμμικό, εκθετικό ή αντίστροφης σιγμοειδούς μορφής~\cite{bengio2015scheduled}. Έτσι το μοντέλο εκπαιδεύεται στις συνθήκες που θα συναντήσει κατά την εκτέλεση και γίνεται ανθεκτικό (\textit{robust}) απέναντι στα σφάλματα που συσσωρεύει κατά την αναδρομική πρόβλεψη.

Στην παρούσα εργασία, το Scheduled Sampling
επεκτείνεται στους αλγόριθμους Gradient Boosting
(LightGBM, XGBoost, Random Forest), ακολουθώντας
ανάλογη λογική: κατά την εκπαίδευση, τα αναδρομικά χαρακτηριστικά υστέρησης αντικαθίστανται με
προβλέψεις του ίδιου του μοντέλου με πιθανότητα $\varepsilon_t$. Εδώ το $\varepsilon$ δηλώνει την πιθανότητα να αντικατασταθεί μια πραγματική υστέρηση από την αντίστοιχη πρόβλεψη του μοντέλου (\textit{sampling probability}), δηλαδή αντίστροφα από την αρχική διατύπωση όπου το $\varepsilon$ ήταν η πιθανότητα χρήσης της πραγματικής τιμής. Ακολουθώντας γραμμικό πρόγραμμα, η $\varepsilon_t$ αυξάνεται από $0{,}10$ σε $0{,}40$ στις τρεις επαναλήψεις
(\textit{SS iterations}), ισορροπώντας τη
διόρθωση προς αναδρομικές συνθήκες με τη
διατήρηση σήματος από τις πραγματικές τιμές.


% ============================================================
\section{Walk-Forward Μηνιαία Επανεκπαίδευση}
\label{sec:walk_forward}
% ============================================================

Η στατική εκπαίδευση ενός μοντέλου σε σταθερό
παράθυρο ιστορικών δεδομένων αδυνατεί να προσαρμοστεί
σε αλλαγές καθεστώτος (\textit{regime shifts})
της αγοράς. Στην ελληνική αγορά, η τιμή
ηλεκτρισμού παρουσίασε ακραία μεταβλητότητα
κατά την περίοδο 2021–2022 (ενεργειακή κρίση)
και σταθεροποίηση το 2024–2025, με αποτέλεσμα
ένα μοντέλο που εκπαιδεύτηκε σε ολόκληρη
αυτή την περίοδο να αποτυγχάνει να αντικατοπτρίσει
τις τρέχουσες συνθήκες.

Η \emph{Walk-Forward Μηνιαία Επανεκπαίδευση}
λύνει αυτό το πρόβλημα με ένα επεκτεινόμενο
παράθυρο δεδομένων: πριν από κάθε μήνα
δοκιμής $m$, το μοντέλο επανεκπαιδεύεται
σε όλα τα διαθέσιμα ιστορικά δεδομένα έως
και το τέλος του προηγούμενου μήνα:
%
\begin{equation}
  \mathcal{D}_m = \left\{(x_t, y_t) : t \leq T_m\right\},
  \quad T_{m+1} = T_m + \Delta
  \label{eq:walk_forward}
\end{equation}
%
όπου $\Delta$ είναι ο αριθμός των ωρών ενός μήνα
(συνήθως $720$–$744$ ώρες) και $T_m$ ο χρόνος
τερματισμού εκπαίδευσης για τον μήνα $m$.
Σε κάθε κύκλο, το μοντέλο ενσωματώνει τα πιο πρόσφατα μοτίβα χωρίς να «ξεχνά»
την ιστορική πληροφορία. Η προσέγγιση διαφέρει
από το rolling window, όπου τα πιο παλαιά
δεδομένα αφαιρούνται ενώ εδώ το σύνολο εκπαίδευσης
διαρκώς αυξάνεται.

Η επανεκπαίδευση ενσωματώνει άμεσα στο μοντέλο τις πρόσφατες μεταβολές στα μοτίβα τιμής και φορτίου, μειώνοντας σημαντικά το σφάλμα σε περιόδους αλλαγής καθεστώτος.
Στο πλαίσιο της παρούσας εργασίας,
η Walk-Forward Μηνιαία Επανεκπαίδευση αξιολογείται
για τρεις μήνες (Δεκέμβριος 2025, Ιανουάριος 2026,
Φεβρουάριος 2026) και συγκρίνεται με τα αντίστοιχα
στατικά μοντέλα, καταγράφοντας βελτίωση
έως~$30\%$ στο MAE.

% ============================================================
% Ολόκληρη η ενότητα (επικεφαλίδα + εισαγωγή + πλατύς πίνακας) τοποθετείται σε μία
% landscape σελίδα. Έτσι αποφεύγεται η κενή portrait σελίδα που γεννούσε το
% \clearpage του landscape όταν η προηγούμενη σελίδα ήταν γεμάτη.
\begin{landscape}
\thispagestyle{empty}%           ← αφαιρεί το header που εμφανίζεται στο πλάι
\noindent\hfill\thepage          % αριθμός σελίδας πάνω-δεξιά στο landscape
\par\vspace{2pt}

\section{Σύνθεση και Τοποθέτηση της Παρούσας Εργασίας}
\label{sec:lit_positioning}

Ο Πίνακας~\ref{tab:related_comparison} συνοψίζει αντιπροσωπευτικές εργασίες πρόβλεψης
τιμής [Electricity Price Forecasting (EPF)] και φορτίου [Electricity Load Forecasting
(ELF)], με την περιοχή μελέτης, τους αλγόριθμους, τον ορίζοντα και τη στρατηγική
πρόβλεψης κάθε εργασίας, και τοποθετεί την παρούσα εργασία ως προς αυτές (τελευταία
γραμμή). Η ποσοτική αντιπαραβολή των ευρημάτων αναπτύσσεται στο
Κεφάλαιο~\ref{ch:experiments}. Κάθε αριθμητική σύγκριση MAE μεταξύ εργασιών είναι
ενδεικτική και όχι αυστηρή, καθώς οι μελέτες διαφέρουν ως προς την περίοδο δοκιμής,
την αγορά, το σύνολο χαρακτηριστικών και τον ορισμό σφάλματος.

\vspace{2pt}
\begin{center}
  \scriptsize
  \setlength{\tabcolsep}{5pt}
  \renewcommand{\arraystretch}{1.1}
  \begin{tabular}{p{3.0cm}p{1.5cm}p{3.4cm}p{6.2cm}p{2.6cm}p{5.8cm}}
    \toprule
    \textbf{Συγγραφέας (Έτος)} & \textbf{Στόχος} & \textbf{Περιοχή (Dataset)} & \textbf{Μέθοδοι} & \textbf{Ορίζοντας} & \textbf{Στρατηγική} \\
    \midrule
    \textit{Kontraros} (2023)~\cite{kontraros2023}
      & Τιμή & Ελλάδα (ΕΧΕ)
      & Ridge, Lasso, Random Forest, LSTM, Quantile Regression Forest
      & Day-Ahead (24h) & Ντετερμινιστική \& πιθανοτική \\
    \addlinespace[4pt]
    \textit{Kanousis} (2025)~\cite{kanousis2025}
      & Τιμή & Ελλάδα (ΕΧΕ / ENTSO-E)
      & SARIMAX, Random Forest, XGBoost, LSTM
      & Day-Ahead (24h) & Rolling-origin (TF, Rec, MIMO, Direct) \\
    \addlinespace[4pt]
    \textit{Μαγκλαράς} (2018)~\cite{magklaras2018}
      & Φορτίο & Ελλάδα (ΔΕΔΔΗΕ) \& Ιταλία (Trento)
      & LSTM, MLP, CNN, ANN
      & VSTLF / STLF & Recursive, Direct, MIMO, υβριδική \\
    \addlinespace[4pt]
    \textit{Michalakopoulos et al.} (2025)~\cite{ntua2025shortwindow}
      & Τιμή & Ελλάδα, Βέλγιο, Ιρλανδία
      & LSTM (FFEC), XGBoost, LightGBM, CatBoost
      & Day-Ahead (24h) & Short training windows (point) \\
    \addlinespace[4pt]
    \textit{Zelios et al.} (2025)~\cite{zelios2025shortterm}
      & Φορτίο & Ελλάδα
      & RNN, LSTM, GRU
      & 1h, 2h, 24h & Μεταβλητό μήκος ακολουθίας (point) \\
    \addlinespace[4pt]
    \textit{Shiblee \& Koukaras} (2025)~\cite{shiblee2025shortterm}
      & Φορτίο & Ελλάδα (δίκτυο διανομής)
      & CNN, LSTM, GRU, CNN-LSTM, LightGBM
      & STLF (24h) & Επιβλεπόμενη μάθηση (point) \\
    \addlinespace[4pt]
    \textit{Lago et al.} (2021)~\cite{lago2021}
      & Τιμή & Ευρώπη (Nord Pool, EPEX, PJM)
      & Deep Learning (DNN), LEAR
      & Day-Ahead (24h) & Ολισθαίνον παράθυρο \\
    \addlinespace[4pt]
    \textit{Hirsch et al.} (2024/2025)
      & Τιμή & Ισπανία, Γερμανία \& Λουξεμβούργο
      & RLin, MLP, FLin, LEAR, DNN, Ensemble (BOA)
      & Day-Ahead (24h) & Ολισθαίνον παράθυρο (Ensemble) \\
    \addlinespace[4pt]
    \textit{Kızıldağ et al.} (2024)~\cite{kizildag2024development}
      & Τιμή & Τουρκία
      & Στατιστικές \& AI μέθοδοι
      & Day-Ahead (24h) & Στατιστική ανάλυση \& AI (point) \\
    \addlinespace[4pt]
    \textit{Narajewski \& Ziel} (2022)
      & Τιμή & Ευρωπαϊκές αγορές
      & Probabilistic DNN, LEAR, Ensemble
      & Day-Ahead \& intraday & Πιθανοτική (rolling, Ensemble) \\
    \addlinespace[4pt]
    \textit{Nikolaidis} (2024)
      & Φορτίο (res.) & Κύπρος (CERA)
      & Νευρωνικά δίκτυα (Levenberg-Marquardt)
      & Day-Ahead (24h) & MIMO \\
    \addlinespace[4pt]
    \midrule
    \textbf{Παρούσα εργασία}
      & Τιμή \& φορτίο & Ελλάδα (ΕΧΕ / ΑΔΜΗΕ)
      & LightGBM, XGBoost, RF, MLP, SVR, Ensembles
      & Day-Ahead (24h) & TF, Rec, MIMO, Direct, Walk-Forward MR \\
    \bottomrule
  \end{tabular}
  % β) λεζάντα κάτω από τον πίνακα
  \captionof{table}{Αντιπροσωπευτικές εργασίες EPF/ELF και τοποθέτηση της παρούσας εργασίας.}
  \label{tab:related_comparison}
\end{center}
\end{landscape}
